\chapter{Introduction}
\pagenumbering{arabic} \setcounter{page}{1}

\section{Chapter Overview}

A number of pages about the background of the project.


\section{Problem Statement}

Provide a list of all the chapters within the thesis and a brief summary of the content.


\section{Research Motivation}

The increasing complexity of global supply chains and consumer demand volatility have intensified the need for intelligent, adaptive, and transparent analytics systems in the consumer goods sector. Recent studies reveal that AI-driven forecasting can reduce prediction errors by 15–20\%, resulting in improved inventory turnover and cost efficiency \autocite{khedr_enhancing_2024}. Nevertheless, all those advantages do not always mean that the introduction of AI solutions will be welcomed by organizations because such solutions are not explainable, do not align well with business goals, and they do not fit the cost structure. The conventional forecasting mechanisms consider over and under-forecasting as an equal offense and ignore their unique financial and service-level consequences \autocite{ledro_artificial_2025}. In addition, the current sustainability requirements and the trends of the circular economy demand reverse logistics returns, recycling, and remanufacturing modeling that are still mostly nonexistent in the majority of analytics models. The rationale behind this research is the necessity to develop a hybrid, interpretable, business-oriented AI that will theorize the gap between data science and strategic supply chain management. This research will contribute to generating practical findings, which will not only provide a better forecast but also a higher level of resilience, profitability, and sustainability in consumer goods supply chains by combining causal explainability, extraneous data streams, and cost-sensitive optimization.

\section{Research Gap}


\section{Research Questions}

This research is informed by a research question which is central and has a series of secondary questions which examine various facets of the problem. The questions below are designed to guide the research and offer an orderly manner of answering research objectives.

\subsubsection{Primary Research Question}

\begin{itemize}
    \item \textbf{RQ1:} How can a hybrid AI-driven analytics framework enhance the accuracy, transparency, and cost-effectiveness of supply chain forecasting in consumer goods companies? 
\end{itemize}

\subsubsection{Secondary Research Questions}

\begin{itemize}
    \item \textbf{RQ2:} What combination of internal and external data sources (e.g., sales history, macroeconomic indicators, social sentiment) most effectively improves forecast performance in consumer goods supply chains? 
    \item \textbf{RQ3:} How can explainable AI (XAI) techniques be integrated into forecasting models to enhance planner trust and decision transparency?
    \item \textbf{RQ4:} In what ways can cost-sensitive and asymmetric optimization functions better align AI-based forecasts with business performance goals such as service level, inventory turnover, and sustainability outcomes?
\end{itemize}

\section{Research Aim and Objectives}

The objectives of the research are as follows and outline the desired direction and quantifiable goals of the research. They will act as a guide on how the research will be carried out and provide some form of consistency on the purpose and the anticipated outcomes of the research.

\subsubsection{Primary Objective}

\begin{itemize}
    \item \textbf{RO1:} To design and evaluate a hybrid AI-driven supply chain analytics framework that integrates interpretability, real-time external data, and cost-sensitive optimization to enhance forecast accuracy and operational resilience in consumer goods companies.
\end{itemize}

\subsubsection{Secondary Objectives}

\begin{itemize}
    \item \textbf{RO2:} To identify and analyze the impact of integrating diverse data sources internal (sales, logistics) and external (economic, sentiment, weather) on demand accuracy. 
    \item \textbf{RO3:} To develop and test explainability mechanisms (e.g., SHAP, LIME) that make AI-based forecasting models transparent and trustworthy for supply chain planners.
    \item \textbf{RO4:} To incorporate asymmetric cost-aware optimization functions within forecasting and replenishment policies, aligning predictive insights with business outcomes such as reduced stockouts, improved inventory turnover, and enhanced sustainability.
\end{itemize}

\section{Significance of the Research}


\section{Scope of the Research}